% Bewertungspaket zum Gesamttest «Darstellungen wechseln» — Quelle des PDFs (python3 scripts/build-lp-pdf.py darstellungen).
% Ganze Punkte je Teilschritt; Werte und Fehlerwerte mit python3 nachgerechnet (scripts/lp/darstellungen/zahlen.py).
% Gerüst aus dem Bewertungspaket von LP 3 (gleichungen-formeln): Auftrag an die KI, Folgefehler, (E)-Zeilen.
\documentclass[11pt]{article}
\usepackage{../lp-druck}
\usepackage{ds-gitter}
\renewcommand{\lpname}{Darstellungen wechseln}
\renewcommand{\lpteilgebiet}{Funktionen ·}
\renewcommand{\lpfuss}{Bewertungspaket}
\newcolumntype{P}{>{\centering\arraybackslash}p{10mm}}
\newenvironment{raster}{\par\smallskip\noindent\renewcommand{\arraystretch}{1.6}\linespread{1.35}\selectfont\tabularx{\linewidth}{@{}>{\raggedright\arraybackslash}p{38mm}P X@{}}\toprule
  \small\textsc{Teilschritt} & \small\textsc{P} & \small\textsc{Musterlösung}\\\midrule}{\bottomrule\endtabularx\par}
\newcommand{\fehler}[1]{\par\smallskip{\small\textcolor{lprot}{\bfseries Typische Fehler:} #1}}
\newcommand{\E}{\,{\footnotesize\bfseries\color{lpgruen}(E)}}
\begin{document}
\lpkopf{Bewertungspaket zum Gesamttest}{}

\begin{lpachtung}{Erst lösen, dann öffnen}
Dieses Paket enthält die Musterlösung. Wer es vor dem Lösen liest, misst am Ende nur, wie gut das Abschreiben gelingt.
\end{lpachtung}

\section*{1 · So gehst du vor}
\begin{enumerate}[leftmargin=*,itemsep=2pt]
\item \textbf{Gesamttest lösen} — vollständig, mit Rechenweg, ohne dieses Paket und ohne Taschenrechner, mit Lineal.
\item \textbf{Selbst bewerten} mit Abschnitt~3: Zeile für Zeile vergleichen und die Punkte eintragen. Das geht ohne Hilfsmittel
  und ist der Normalfall.
\item \textbf{Oder bewerten lassen.} Ob du dafür eine KI nutzt und welche, entscheidest du selbst und in eigener Verantwortung —
  je nachdem, was dir aufgrund deines Alters und deiner persönlichen Situation erlaubt ist (Altersgrenzen und
  Nutzungsbedingungen des Anbieters, allenfalls Einverständnis deiner Eltern). Fotografiere oder scanne deine Lösung gut lesbar,
  \textbf{ohne Namen}. Handyfotos tragen oft unsichtbar Standort, Zeit und Gerät mit (Metadaten) — lade lieber einen Scan oder ein
  Bildschirmfoto des Fotos hoch. Lade nur deine Lösung, das Testblatt und dieses PDF hoch, keine persönlichen Daten, und füge den Auftrag aus
  Abschnitt~2 ein.
\item \textbf{Rückmeldung prüfen.} Stimmt, was die KI aus deiner Handschrift gelesen hat? Eine KI kann sich irren —
  im Zweifel gilt das Raster in Abschnitt~3.
\item \textbf{Gezielt wiederholen} nach Abschnitt~4.
\end{enumerate}

\section*{2 · Auftrag an die KI}
\begin{tcolorbox}[colback=lplinie!25,colframe=lplinie,boxrule=0.5pt,arc=1mm,fontupper=\small]
Du bist Korrektorin oder Korrektor für Mathematik auf der Sekundarstufe I in der Schweiz. Im Anhang findest du (1)~das Testblatt mit den Aufgaben, (2)~das Bewertungspaket zum Gesamttest «Darstellungen wechseln» mit Musterlösung und Punkteraster und (3)~Fotos meiner handschriftlichen Lösung.\par\medskip
Geh so vor:\par
1. Schreib zu jeder Aufgabe G1 bis G8 zuerst ab, was du in meiner Lösung liest — Rechenweg, Ergebnis und Begründung; bei der Zeichnung (G5), wie meine Achsen beschriftet und eingeteilt sind und wo meine Strecke beginnt und endet. Was du nicht sicher lesen kannst, markierst du mit [unsicher]. Nicht raten. Achte besonders auf Minuszeichen und Dezimalpunkte.\par
2. Vergib die Punkte genau nach dem Raster im Bewertungspaket (Abschnitt 3), ganze Punkte je Zeile. Ziehe einen Fehler nur einmal ab: Rechne ich mit einem falschen Zwischenergebnis richtig weiter, gibt es die Folgepunkte. Ausnahme: Zeilen mit \textbf{(E)} gibt es nur für das richtige Ergebnis, nie als Folgepunkt. Die Zeilen «Typische Fehler» sagen, welche Rasterzeile bei diesem Fehler 0 Punkte gibt — sie sind kein zusätzlicher Abzug.\par
3. Andere richtige Lösungswege zählen voll, zum Beispiel ein anderes Wertepaar zum Einsetzen oder eine andere sinnvolle Einteilung der Achsen. Antworten in Worten zählen, wenn sie dasselbe sagen wie die Musterlösung (Grösse und Zahl mit Einheit). Gleichwertig: $y = 0.3x + 1.5$ und $y = \tfrac{3}{10}x + \tfrac{3}{2}$; $y = 18 - x$ und $y = -x + 18$; Brüche und Dezimalzahlen mit gleichem Wert. Eine Begründung «stimmt nicht» zählt nur mit einem Zahlenbeispiel oder einer Rechnung.\par
4. Begründe jeden Abzug in einem Satz und nenne die Fehlerart (zum Beispiel «Schritt in x vergessen», «nicht durch 2 geteilt», «Skala ungleichmässig», «ganze Gerade statt Strecke», «linear fallend mit antiproportional verwechselt», «Beispiel statt Gegenbeispiel»). Schreib die Musterlösung nicht ab — zeig mir, wo mein Fehler liegt.\par
5. Gib zum Schluss eine Tabelle «Aufgabe | Punkte | Begründung», die Gesamtpunktzahl (von 25) und — nach der Selbsteinschätzung in Abschnitt 4 — welches Kapitel des Leitprogramms ich wiederholen soll. Keine Note.\par
6. Fehlt ein Foto oder ist eine Aufgabe unlesbar, sag es, statt sie zu bewerten.
\end{tcolorbox}

\needspace{16\baselineskip}
\section*{3 · Musterlösung und Punkteraster}
\textbf{Allgemein:} Ganze Punkte je Zeile. Ein Fehler wird nur einmal abgezogen (Folgefehler). Gleichwertige Wege zählen voll.
In der Spalte P: \textbf{(E)}~= Ergebnispunkt, Ergebnis richtig \emph{und} Weg erkennbar; ohne Zeichen = die
Zeile bewertet den Weg, die Zeichnung oder die Begründung. Der ganze Test ist ohne Taschenrechner gestellt.
«Typische Fehler» nennt, welche Zeile 0 Punkte gibt, und die Punkte, die bleiben.
\textbf{Teilrichtige Zeilen:} Verlangt eine Zeile mehrere Angaben und stimmt nur ein Teil, gibt die Zeile 0 Punkte,
ausser die Zeile nennt ausdrücklich eine Aufteilung.
\textbf{Zahlen:} Dezimalzahl oder Bruch, wenn gleich ($7.5 = \tfrac{15}{2}$). Einheiten sind in Rechnungen nicht verlangt, in Antworten in Worten schon.

\begin{lpaufgabe}{G1}{E-Scooter}{4 P}
\begin{raster}
(a) Gleichung & 1 & $y = 0.3x + 1.5$ ($1.50$ Fr. einmal, $0.30$ Fr. pro Minute)\\
(b) Tabelle & 1 & $x = 0, 10, 20$: $y = 1.5,\ 4.5,\ 7.5$. Auch als Folgepunkt mit der eigenen Gleichung.\\
(c) Punkt & 1 & Eine Fahrt von $20$ Minuten kostet insgesamt $7.50$ Fr. (Zeit und Gesamtpreis müssen beide vorkommen.)\\
(d) & 1\E & $0.3x + 1.5 = 9$, $0.3x = 7.5$, $x = 25$ Minuten\\
\end{raster}
\fehler{$m$ und $b$ vertauscht ($y = 1.5x + 0.3$): Zeile (a) 0; (b) als Folgepunkt ($0.3$, $15.3$, $30.3$) 1.
(c) «Eine Minute kostet $7.50$ Fr.» oder «$7.50$ Minuten kosten $20$ Fr.»: Zeile (c) 0.
(d) $9 : 0.3 = 30$ (Entsperren vergessen): Zeile (d) 0.}
\end{lpaufgabe}

\begin{lpaufgabe}{G2}{Tabelle}{3 P}
\begin{raster}
$m$ & 1 & $m = \dfrac{25 - 17}{5 - 3} = \dfrac{8}{2} = 4$ (oder mit $9$: $\tfrac{16}{4} = 4$)\\
$b$ & 1 & $17 = 4 \cdot 3 + b$, $b = 5$. Auch als Folgepunkt mit dem eigenen $m$.\\
Gleichung & 1\E & $y = 4x + 5$ (Probe mit $x = 9$: $41$)\\
\end{raster}
\fehler{Schritt in $x$ vergessen ($m = 8$): $m$-Zeile und Gleichungszeile 0, $b$-Zeile als Folgepunkt ($b = -7$) 1 $\to$ 1 von 3\,P.
$b = 17$ (erster $y$-Wert): $b$-Zeile und Gleichungszeile 0 $\to$ 1 von 3\,P.}
\end{lpaufgabe}

\begin{lpaufgabe}{G3}{Rechteck}{3 P}
\begin{raster}
(a) Ansatz & 1 & $2x + 2y = 36$, durch $2$: $x + y = 18$\\
(a) Ergebnis & 1\E & $y = 18 - x$\\
(b) & 1 & $y = 18 - 7.5 = 10.5$ cm. Auch als Folgepunkt mit der eigenen Gleichung, wenn $y > 0$.\\
\end{raster}
\fehler{Nicht durch $2$ geteilt ($y = 36 - 2x$): Ergebniszeile (a) 0; (b) als Folgepunkt ($21$) 1 $\to$ 2 von 3\,P.
$y = 36 - x$ (jede Seite nur einmal gezählt): Ansatz- und Ergebniszeile 0; (b) als Folgepunkt ($28.5$) 1 $\to$ 1 von 3\,P.}
\end{lpaufgabe}

\begin{lpaufgabe}{G4}{Kino-Gutschein}{3 P}
\begin{raster}
(a) & 1 & $8y = 60 - 12x$, $y = 7.5 - 1.5x$ (gleichwertig $y = -\tfrac{3}{2}x + \tfrac{15}{2}$)\\
(b) & 1 & Für jeden Erwachsenen mehr gibt es $1.5$ Kinderbillette weniger (oder: $2$ Erwachsene mehr, $3$ Kinder weniger). Auch als Folgepunkt passend zur eigenen Steigung.\\
(c) & 1\E & Nur ganze Anzahlen: $(1 \mid 6)$, $(3 \mid 3)$, $(5 \mid 0)$ — alle drei, keine falschen ($(0 \mid 7.5)$ ist nicht möglich). Ob $(5 \mid 0)$ «Kinder» braucht, ist egal: Es zählt auch ohne.\\
\end{raster}
\fehler{(a) Nicht durch $8$ geteilt ($y = 60 - 12x$): Zeile (a) 0.
(b) «Ein Kind kostet $1.5$ Fr.» oder Steigung ohne Deutung: Zeile (b) 0.
(c) $(0 \mid 7.5)$ aufgeführt oder ein Paar fehlt: Zeile (c) 0.}
\end{lpaufgabe}

\begin{lpaufgabe}{G5}{Pool}{4 P}
\begin{center}
\leernetz{12}{9}{\draw[->,line width=0.6pt] (0,0) -- (12.5,0) node[right] {$x$ [h]};
  \draw[->,line width=0.6pt] (0,0) -- (0,9.5) node[above] {$y$ [cm]};
  \foreach \i in {2,4,...,12} \node[below] at (\i,0) {$\i$};
  \foreach \j/\w in {2/40,4/80,6/120,8/160} \node[left] at (0,\j) {$\w$};
  \draw[lpblau,line width=1.2pt] (0,2) -- (8,8);
  \fill[lporange] (0,2) circle (2.4pt); \fill[lpgruen] (8,8) circle (2.4pt);}
\end{center}
\begin{raster}
(a) & 1 & $y = 15x + 40$; voll bei $15x + 40 = 160$, $x = 8$: $0 \le x \le 8$ Stunden\\
(b) & 1 & $40 \le y \le 160$ cm. Auch als Folgepunkt mit dem eigenen Bereich aus (a).\\
(c) Achsen & 1 & Beide Achsen mit Grösse und Einheit ($x$ in h, $y$ in cm) und gleichmässiger Einteilung, die bis $8$ h und $160$ cm reicht (z.\,B.\ $1$ Häuschen $= 1$ h und $1$ Häuschen $= 20$ cm wie im Bild).\\
(d) Strecke & 1 & Strecke von $(0 \mid 40)$ bis $(8 \mid 160)$, nicht darüber hinaus. Mit eigenen Endpunkten aus (a) als Folgepunkt.\\
\end{raster}
\fehler{Ganze Gerade über $x = 8$ hinaus: Zeile (d) 0.
Achse ohne Einheit oder mit ungleichmässiger Einteilung: Zeile (c) 0.
Bereich $0 \le x \le 160$ oder $0 \le x \le 10.\overline{6}$ ($160 : 15$, Anfangswasser vergessen): Zeile (a) 0; (d) als Folgepunkt nur mit passender Strecke.}
\end{lpaufgabe}

\begin{lpaufgabe}{G6}{Achse}{2 P}
\begin{raster}
(a) & 1 & Die Abstände sind gleich, stehen aber für verschiedene Schritte ($10$, $10$, $30$, $50$): Die Skala ist nicht gleichmässig.\\
(b) & 1 & Bis $20$ steht ein Häuschen für $10$, danach für $30$ und $50$: Die Punkte rücken oben zusammen, die Strecken werden flacher. Mit gleichmässiger Skala lägen alle Punkte auf einer Geraden.\\
\end{raster}
\fehler{«Die $30$ und $40$ fehlen»: Das trifft den Kern nur, wenn gesagt wird, dass die Abstände nicht passen; sonst Zeile (a) 0.
(b) «Weil die Funktion nicht linear ist»: Zeile (b) 0.}
\end{lpaufgabe}

\begin{lpaufgabe}{G7}{Zwei Tabellen}{3 P}
\begin{raster}
A & 1 & Antiproportional: $1 \cdot 30 = 2 \cdot 15 = 3 \cdot 10 = 6 \cdot 5 = 30$, das Produkt bleibt gleich.\\
B & 1 & Linear fallend: Pro Schritt um $1$ in $x$ wird $y$ um $5$ kleiner ($25, 20, 15$; von $3$ bis $6$: $-15$ auf $3$ Schritte).\\
Gleichung & 1\E & B: $y = -5x + 30$ (Probe $x = 6$: $0$)\\
\end{raster}
\fehler{A «linear fallend», weil $y$ kleiner wird: Zeile A 0.
B «antiproportional» ohne Rechnung: Zeile B 0.
Nur Behauptung ohne Rechnung bei A oder B: diese Zeile 0.}
\end{lpaufgabe}

\begin{lpaufgabe}{G8}{Aussagen prüfen}{3 P}
\begin{raster}
(a) Rechnung & 1 & $x + y = 10$: $x = 2$, $y = 8$; $x = 4$, $y = 6$\\
(a) Urteil & 1 & Falsch: $6$ ist nicht die Hälfte von $8$. Auch als Folgepunkt passend zur eigenen Rechnung.\\
(b) & 1 & Falsch, mit Gegenbeispiel: $2 \cdot 8 = 16$ und $4 \cdot 6 = 24$ cm² (jedes Paar mit verschiedenen Flächen zählt; ein gedrehtes Rechteck wie $2 \times 8$ und $8 \times 2$ nicht)\\
\end{raster}
\fehler{(a) $y = 20 - x$ (nicht durch $2$ geteilt): Rechnungszeile 0; Urteil als Folgepunkt 1.
(b) Nur «stimmt nicht» ohne Zahlen, oder ein einzelnes Beispiel, das passt: Zeile (b) 0.}
\end{lpaufgabe}

\needspace{12\baselineskip}
\section*{4 · Selbsteinschätzung}
\begin{tabularx}{\linewidth}{@{}l X@{}}
\toprule
\textbf{22 – 25 P} & Die geprüften Teile sitzen. Wo du Punkte verloren hast: das Kapitel dieser Aufgabe nochmals.\\
\textbf{17 – 21 P} & Den schwächsten Teil nochmals, gemessen am Anteil (Teil A 7, B 6, C 6, D 6 P): Tüfteln und Übungen des Kapitels dieses Teils.\\
\textbf{11 – 16 P} & Zurück zu den Kapiteln aller Aufgaben, in denen du Punkte verloren hast.\\
\textbf{0 – 10 P} & Zurück zu den Kapiteln jedes Teils, in dem du weniger als die Hälfte erreicht hast; sind es alle vier, zum Vortest und zu Kapitel 1.\\
\bottomrule
\end{tabularx}

\smallskip
\textbf{Aufgabe $\to$ Kapitel:} G1, G2 $\to$ 1; G3, G4 $\to$ 2; G5, G6 $\to$ 3; G7, G8 $\to$ 4.
\end{document}
